% !TEX root = ../main.tex
%% Corrigés — Chapitre 1 — Suites numériques
%% Tous les résultats ont été vérifiés par calcul formel (sympy).

\corrigesChapitre{1}{Suites numériques}

\begin{deuxcolonnes}

\begin{corrige}{1}{Limites usuelles}
\begin{questions}
  \item Limite de référence : $\res{\lim u_n=0}$
  \item Limite de référence : $\res{\lim v_n=0}$
  \item Limite de référence : $\res{\lim w_n=+\infty}$
  \item Limite de référence : $\res{\lim t_n=+\infty}$
  \item $\lim\frac{2}{n}=0$, donc $\res{\lim s_n=5}$
\end{questions}
\end{corrige}

\begin{corrige}{2}{Opérations sur les limites}
\begin{questions}
  \item $\lim n=+\infty$, donc $\res{\lim(n+7)=+\infty}$
  \item \methode{simplifier d'abord}\\
    Pour $n\geq1$, $3n\times\frac{1}{n}=3$ : la limite vaut $\res{3}$. Sans simplifier, on obtiendrait la forme indéterminée $\infty\times0$.
  \item $\frac{n+1}{n}=1+\frac{1}{n}$, donc la limite vaut $\res{1}$
  \item $\lim\frac{5}{n^{2}}=0$, donc la limite vaut $\res{2}$
\end{questions}
\end{corrige}

\begin{corrige}{3}{Suites géométriques}
\begin{questions}
  \item $u_n=2\times3^{n}$ avec $3>1$ et $2>0$ :\\ $\res{\lim u_n=+\infty}$
  \item $v_n=5\times\left(\frac12\right)^{n}$ avec $-1<\frac12<1$ :\\ $\res{\lim v_n=0}$
  \item $-1<\frac13<1$, donc $\res{\lim w_n=0}$
  \item $q=-2\leq-1$ : $\res{(t_n)\text{ n'a pas de limite}}$\\
    En effet, $t_{2p}=4^{p}$ tend vers $+\infty$ alors que $t_{2p+1}=-2\times4^{p}$ tend vers $-\infty$.
\end{questions}
\end{corrige}

\begin{corrige}{4}{Formes indéterminées}
\methode{factoriser par le terme dominant}
\begin{questions}
  \item $n^{2}-n=n^{2}\left(1-\frac{1}{n}\right)$ : $\res{+\infty}$
  \item $\dfrac{n^{2}\left(3+\frac{2}{n}+\frac{1}{n^{2}}\right)}{n^{2}\left(5-\frac{4}{n^{2}}\right)}$ tend vers $\res{\dfrac{3}{5}}$
  \item $\dfrac{2n^{3}-n+1}{n^{2}+3n}=n\times\dfrac{2-\frac{1}{n^{2}}+\frac{1}{n^{3}}}{1+\frac{3}{n}}$ :\\ $\res{+\infty}$
  \item $n-\sqrt{n}=\sqrt{n}\left(\sqrt{n}-1\right)$ : $\res{+\infty}$
  \item $\dfrac{4n+3}{2n^{3}-n^{2}+1}=\dfrac{1}{n^{2}}\times\dfrac{4+\frac{3}{n}}{2-\frac{1}{n}+\frac{1}{n^{3}}}$ :\\ $\res{0}$
  \item \methode{quantité conjuguée}\\
    $\sqrt{n^{2}+n}-n=\dfrac{n}{\sqrt{n^{2}+n}+n}$\\
    $\phantom{\sqrt{n^{2}+n}-n}=\dfrac{1}{\sqrt{1+\frac{1}{n}}+1}$\\
    qui tend vers $\res{\dfrac{1}{2}}$
  \item Le quotient des termes dominants vaut $\frac{n^{4}}{n^{4}}=1$ : la limite vaut $\res{1}$
\end{questions}
\end{corrige}

\begin{corrige}{5}{Théorème des gendarmes (1)}
\methode{encadrer, puis conclure}
\begin{questions}
  \item $-\frac{1}{n}\leq\frac{\sin n}{n}\leq\frac{1}{n}$ : $\res{0}$
  \item $\dfrac{n^{2}-1}{n^{2}+1}\leq v_n\leq\dfrac{n^{2}+1}{n^{2}+1}=1$, et la borne de gauche tend vers $1$ : $\res{1}$
  \item $-\frac{1}{n}\leq\frac{(-1)^{n}}{n}\leq\frac{1}{n}$ : $\res{0}$
  \item $\dfrac{3n-2}{n+5}\leq v_n\leq\dfrac{3n+2}{n+5}$, les deux bornes tendent vers $3$ : $\res{3}$
  \item Pour $n\geq2$ :\\ $\dfrac{n^{3}-1}{n^{3}-2}\leq v_n\leq\dfrac{n^{3}+1}{n^{3}-2}$ : $\res{1}$
  \item $\dfrac{n-1}{n^{2}}\leq v_n\leq\dfrac{n+1}{n^{2}}$ : $\res{0}$
  \item $\dfrac{5n^{2}-4}{5n^{2}+2}\leq v_n\leq\dfrac{5n^{2}+2}{5n^{2}+2}=1$ : $\res{1}$
\end{questions}
\end{corrige}

\begin{corrige}{6}{Théorème des gendarmes (2)}
\begin{questions}
  \item $-\sqrt{n}\leq(-1)^{n}\sqrt{n}\leq\sqrt{n}$, donc\\
    $\dfrac{n-\sqrt{n}}{n+1}\leq v_n\leq\dfrac{n+\sqrt{n}}{n+1}$.\\
    Or $\dfrac{n\pm\sqrt{n}}{n+1}=\dfrac{1\pm\frac{1}{\sqrt{n}}}{1+\frac{1}{n}}$ tend vers $1$ : $\res{1}$
  \item Comme $2n^{2}-n+1>0$ :\\
    $\dfrac{2n^{2}-n}{2n^{2}-n+1}\leq v_n\leq\dfrac{2n^{2}+n}{2n^{2}-n+1}$ : $\res{1}$
  \item $-2\leq\sin n+\cos n\leq2$, donc\\ $-\frac{2}{n}\leq v_n\leq\frac{2}{n}$ : $\res{0}$
  \item $\dfrac{n^{4}-n^{2}}{n^{4}+3n}\leq v_n\leq\dfrac{n^{4}+n^{2}}{n^{4}+3n}$ : $\res{1}$
  \item $\dfrac{4n-3}{4n+3}\leq v_n\leq\dfrac{4n+1}{4n+3}$ : $\res{1}$
  \item $\dfrac{n^{2}-n-1}{n^{2}+2}\leq v_n\leq\dfrac{n^{2}-n+1}{n^{2}+2}$ : $\res{1}$
\end{questions}
\end{corrige}

\begin{corrige}{7}{Récurrence et sommes}
\begin{questions}
  \item \etape{Initialisation} $1=\frac{1\times2}{2}$.\\
    \etape{Hérédité} Si $1+\dots+k=\frac{k(k+1)}{2}$, alors\\
    $1+\dots+k+(k+1)=\frac{k(k+1)}{2}+(k+1)$\\
    $=\frac{(k+1)(k+2)}{2}$.\\
    \etape{Conclusion} $\res{\text{vraie pour tout }n\geq1}$
  \item \etape{Initialisation} $1^{2}=\frac{1\times2\times3}{6}$.\\
    \etape{Hérédité} $\frac{k(k+1)(2k+1)}{6}+(k+1)^{2}$\\
    $=\frac{(k+1)\left(2k^{2}+7k+6\right)}{6}$\\
    $=\frac{(k+1)(k+2)(2k+3)}{6}$,\\
    ce qui est la formule au rang $k+1$.
  \item \etape{Initialisation} $2^{0}=1=2^{1}-1$.\\
    \etape{Hérédité} $\left(2^{k+1}-1\right)+2^{k+1}=2^{k+2}-1$.
  \item \etape{Initialisation} $\frac{1}{1\times2}=\frac12$.\\
    \etape{Hérédité} $\frac{k}{k+1}+\frac{1}{(k+1)(k+2)}=\frac{k^{2}+2k+1}{(k+1)(k+2)}$\\
    $=\frac{(k+1)^{2}}{(k+1)(k+2)}=\frac{k+1}{k+2}$.
\end{questions}
\end{corrige}

\begin{corrige}{8}{Récurrence, inégalités et divisibilité}
\begin{questions}
  \item \etape{Initialisation} $2^{0}=1\geq1$.\\
    \etape{Hérédité} $2^{k+1}=2\times2^{k}\geq2(k+1)$\\ $=2k+2\geq k+2$.
  \item \etape{Initialisation} $3^{1}=3\geq3$.\\
    \etape{Hérédité} $3^{k+1}\geq3\left(2k^{2}+1\right)=6k^{2}+3$, et
    $6k^{2}+3-\left(2(k+1)^{2}+1\right)=4k(k-1)\geq0$ pour $k\geq1$.
  \item \etape{Initialisation} $4^{0}-1=0=3\times0$.\\
    \etape{Hérédité} Si $4^{k}-1=3m$ avec $m$ entier :\\
    $4^{k+1}-1=4\left(4^{k}-1\right)+3$\\ $=3(4m+1)$.
  \item \etape{Initialisation} $2^{4}=16\geq4^{2}$.\\
    \etape{Hérédité} $2^{k+1}\geq2k^{2}$, et $2k^{2}-(k+1)^{2}=(k-1)^{2}-2>0$ pour $k\geq4$.
\end{questions}
\end{corrige}

\begin{corrige}{9}{Suites arithmético-géométriques}
\methode{point fixe $\ell=\frac{b}{1-a}$, puis $u_n=\ell+(u_0-\ell)a^{n}$}
\begin{questions}
  \item $\ell=-3$ : $u_n=-3+4\times2^{n}$\\
    $\res{u_n=2^{n+2}-3}$, qui tend vers $+\infty$.
  \item $\ell=4$ : $\res{v_n=4+\left(\frac12\right)^{n}}$\\ qui converge vers $4$.
  \item $\ell=3$ : $\res{w_n=3-3^{n+1}}$\\ qui tend vers $-\infty$.
  \item $\ell=20$ : $t_n=20-10\times0{,}8^{n}$\\ converge vers $\res{20}$.
  \item $\ell=\frac32$ : $s_n=\frac32+\frac32\left(-\frac13\right)^{n}$\\
    converge vers $\res{\frac{3}{2}}$ car $-1<-\frac13<1$.
\end{questions}
\end{corrige}

\end{deuxcolonnes}

%% ---------------------------------------------------------------------
\begin{corrige}{10}{Une suite récurrente avec une racine carrée}
\begin{questions}
  \item \etape{Initialisation} $u_0=1$ et $0\leq1\leq2$.\enspace
    \etape{Hérédité} Si $0\leq u_k\leq2$, alors $2\leq u_k+2\leq4$, et la fonction racine carrée
    étant croissante, $\sqrt2\leq u_{k+1}\leq2$, donc $0\leq u_{k+1}\leq2$.\enspace
    \etape{Conclusion} $\res{0\leq u_n\leq2}$ pour tout entier naturel $n$.
  \item \methode{multiplier par la quantité conjuguée}\\
    $u_{n+1}-u_n=\sqrt{u_n+2}-u_n=\dfrac{u_n+2-u_n^{2}}{\sqrt{u_n+2}+u_n}=\dfrac{(2-u_n)(1+u_n)}{\sqrt{u_n+2}+u_n}$.\\
    D'après la question 1, $2-u_n\geq0$, $1+u_n>0$ et le dénominateur est strictement positif :
    $u_{n+1}-u_n\geq0$, donc $\res{(u_n)\text{ est croissante}}$.
  \item La suite $(u_n)$ est croissante et majorée par $2$. D'après le théorème de convergence
    monotone, elle $\res{\text{converge}}$.
  \item $\ell=\sqrt{\ell+2}$ entraîne $\ell^{2}=\ell+2$, soit $\ell^{2}-\ell-2=(\ell-2)(\ell+1)=0$.
    Comme $u_n\geq0$ pour tout $n$, $\ell\geq0$, donc $\res{\ell=2}$. On vérifie : $\sqrt{2+2}=2$.
\end{questions}
\end{corrige}

%% ---------------------------------------------------------------------
\begin{corrige}{11}{Devoir de synthèse}
\begin{deuxcolonnes}
\partieExo{A}{Limites et opérations}
\begin{questions}
  \item $\dfrac{2+\frac{3}{n}-\frac{1}{n^{2}}}{1+\frac{5}{n^{2}}}$ tend vers $\res{2}$
  \item $n^{3}\left(1-\frac{2}{n}\right)$ : $\res{+\infty}$
  \item $0<\frac23<1$, donc $\res{\lim u_n=0}$
\end{questions}
\partieExo{B}{Théorème des gendarmes}
\begin{questions}[resume]
  \item $\dfrac{2n-1}{n+1}\leq v_n\leq\dfrac{2n+1}{n+1}$ : $\res{2}$
  \item Pour $n\geq2$, $0<n^{2}-1\leq n^{2}+\cos n\leq n^{2}+1$, donc $\dfrac{n^{2}-1}{n^{2}+1}\leq v_n\leq1$ : $\res{1}$
\end{questions}
\columnbreak
\partieExo{C}{Une suite récurrente}
\begin{questions}[resume]
  \item \etape{Initialisation} $v_0=0$.\\
    \etape{Hérédité} Si $0\leq v_k\leq1$, alors $3\leq v_k+3\leq4$, donc $\frac34\leq v_{k+1}\leq1$.\\
    $\res{0\leq v_n\leq1}$
  \item $v_{n+1}-v_n=\dfrac{v_n+3}{4}-v_n=\dfrac{3(1-v_n)}{4}\geq0$ : $(v_n)$ est croissante.
  \item Croissante et majorée par $1$, $(v_n)$ converge vers $\ell$. En passant à la limite dans
    $v_{n+1}=\frac{v_n+3}{4}$ : $\ell=\frac{\ell+3}{4}$, soit $\res{\ell=1}$.
\end{questions}
\end{deuxcolonnes}
\end{corrige}

%% ---------------------------------------------------------------------
\begin{corrige}{12}{Une suite homographique}
\partieExo{A}{Étude de la suite $(u_n)$}
\begin{questions}
  \item \begin{questions}
      \item $u_1=\dfrac{2/3}{4/3}=\res{\dfrac12}$\quad et\quad $u_2=\dfrac{1}{3/2}=\res{\dfrac23\approx0{,}667}$
      \item \etape{Initialisation} $0<\frac13<1$.\enspace
        \etape{Hérédité} Si $0<u_k<1$ : $u_{k+1}>0$ (quotient de deux nombres positifs), et
        $u_{k+1}<1\iff2u_k<1+u_k\iff u_k<1$, ce qui est vrai. Donc $\res{0<u_n<1}$.
    \end{questions}
  \item \begin{questions}
      \item $u_{n+1}-u_n=\dfrac{2u_n-u_n(1+u_n)}{1+u_n}=\dfrac{u_n-u_n^{2}}{1+u_n}=\dfrac{u_n(1-u_n)}{1+u_n}$.
      \item Comme $0<u_n<1$, les trois facteurs sont strictement positifs : $\res{(u_n)\text{ est croissante}}$.
    \end{questions}
\end{questions}
\partieExo{B}{Une suite auxiliaire}
\begin{questions}[resume]
  \item \begin{questions}
      \item $v_{n+1}=\dfrac{u_{n+1}}{1-u_{n+1}}=\dfrac{\frac{2u_n}{1+u_n}}{\frac{1+u_n-2u_n}{1+u_n}}=\dfrac{2u_n}{1-u_n}=2v_n$.
      \item $(v_n)$ est $\res{\text{géométrique de raison }2}$, de premier terme $v_0=\dfrac{1/3}{2/3}=\dfrac12$.
    \end{questions}
  \item \begin{questions}
      \item $v_n=\frac12\times2^{n}=\res{2^{n-1}}$.
      \item De $v_n(1-u_n)=u_n$, on tire $u_n=\dfrac{v_n}{1+v_n}=\dfrac{2^{n-1}}{1+2^{n-1}}=\res{\dfrac{2^{n}}{2^{n}+2}}$.
    \end{questions}
\end{questions}
\partieExo{C}{Limite}
\begin{questions}[resume]
  \item $2>1$ et $v_0>0$ : $\lim v_n=+\infty$. Comme $u_n=\dfrac{1}{1+\frac{1}{v_n}}$ et $\lim\frac{1}{v_n}=0$ : $\res{\lim u_n=1}$.
  \item $\dfrac{2\times1}{1+1}=1$ : $\ell=1$ est solution. Les solutions de $x=f(x)$, avec $f(x)=\dfrac{2x}{1+x}$, sont
    les \emph{points fixes} de $f$ (ici $0$ et $1$) : une suite $u_{n+1}=f(u_n)$ convergente, avec $f$ continue,
    a pour limite un point fixe de $f$.
\end{questions}
\end{corrige}

%% ---------------------------------------------------------------------
\begin{corrige}{13}{Une suite qui converge très vite}
\noindent
\begin{minipage}[t]{0.56\linewidth}\vspace{0pt}\raggedright
\partieExo{A}{Conjecture graphique}
\begin{questions}
  \item $u_1=\frac12\times\frac32=\res{\frac34}$,\quad $u_2=\frac34\times\frac54=\res{\frac{15}{16}}$,\\
    $u_3=\frac{15}{16}\times\frac{17}{16}=\res{\frac{255}{256}}$.
  \item Voir la figure : la suite semble \textbf{croissante} et \textbf{converger vers $1$}, abscisse du point
    d'intersection de $\Cf$ et de la droite $y=x$.
\end{questions}
\partieExo{B}{Convergence}
\begin{questions}[resume]
  \item $f(x)=2x-x^{2}$, donc $f'(x)=2(1-x)$.\\
    Sur $\intFF{0}{1}$, $f'(x)\geq0$ et $f'$ ne s'annule qu'en $1$ :\\
    $\res{f\text{ est strictement croissante sur }\intFF{0}{1}}$\\
    avec $f(0)=0$ et $f(1)=1$.
\end{questions}
\end{minipage}\hfill
\begin{minipage}[t]{0.4\linewidth}\vspace{0pt}
\centering
\begin{tikzpicture}
\begin{axis}[otGraphe, width=58mm, xmin=0, xmax=1.08, ymin=0, ymax=1.08,
  xtick={0.5,1}, ytick={0.5,1}, xticklabels={,}, yticklabels={$\frac12$,$1$},
  minor tick num=4, grid=both, minor grid style={otRule!45,line width=0.2pt}]
  \addplot[otRule!80!otMuted,line width=0.6pt,domain=0:1.05] {x};
  \addplot[otNavy,line width=1.1pt,domain=0:1,samples=80] {x*(2-x)};
  \addplot[otGreen,line width=0.8pt] coordinates {(0.5,0) (0.5,0.75) (0.75,0.75)
    (0.75,0.9375) (0.9375,0.9375) (0.9375,0.99609) (0.99609,0.99609)};
  \draw[otGreen!70!otMuted,densely dotted,line width=0.5pt] (axis cs:0.75,0.75) -- (axis cs:0.75,0);
  \draw[otGreen!70!otMuted,densely dotted,line width=0.5pt] (axis cs:0.9375,0.9375) -- (axis cs:0.9375,0);
  \addplot[only marks,mark=*,mark size=1.3pt,otGreen] coordinates {(0.5,0) (0.75,0) (0.9375,0)};
  \node[anchor=north,font=\scriptsize] at (axis cs:0.5,-0.05) {$u_0$};
  \node[anchor=north,font=\scriptsize] at (axis cs:0.75,-0.05) {$u_1$};
  \node[anchor=north,font=\scriptsize] at (axis cs:0.94,-0.05) {$u_2$};
  \node[anchor=south east,font=\scriptsize,text=otNavy] at (axis cs:0.3,0.62) {$\Cf$};
\end{axis}
\end{tikzpicture}\\[0.5mm]
\legende{$u_3\approx0{,}996$ : les termes se rapprochent de $1$ très vite.}
\end{minipage}

\begin{questions}[start=4]
  \item \etape{Initialisation} $0<\frac12<\frac34<1$.\enspace
    \etape{Hérédité} Si $0<u_k<u_{k+1}<1$, comme $f$ est strictement croissante sur $\intFF{0}{1}$ :
    $f(0)<f(u_k)<f(u_{k+1})<f(1)$, c'est-à-dire $0<u_{k+1}<u_{k+2}<1$.\enspace
    \etape{Conclusion} $\res{0<u_n<u_{n+1}<1}$ pour tout entier naturel $n$.
  \item $(u_n)$ est croissante et majorée par $1$ : elle converge.
\end{questions}
\partieExo{C}{Expression explicite et vitesse de convergence}
\begin{questions}[start=6]
  \item $v_{n+1}=1-u_{n+1}=1-2u_n+u_n^{2}=(1-u_n)^{2}=\res{v_n^{2}}$.
  \item \etape{Initialisation} $v_0=\frac12=\left(\frac12\right)^{2^{0}}$.\enspace
    \etape{Hérédité} $v_{k+1}=v_k^{2}=\left(\left(\frac12\right)^{2^{k}}\right)^{2}=\left(\frac12\right)^{2\times2^{k}}=\res{\left(\frac12\right)^{2^{k+1}}}$.
  \item Comme $2^{n}\geq n$ et $0<\frac12<1$ : $0<v_n\leq\left(\frac12\right)^{n}$. Or $\lim\left(\frac12\right)^{n}=0$, donc,
    d'après le théorème des gendarmes, $\lim v_n=0$, et $\res{\lim u_n=\lim(1-v_n)=1}$.
  \item Les valeurs successives de \lstinline[style=otPython]!v! sont $0{,}25$ ; $0{,}0625$ ; $0{,}0039$ ;
    $1{,}5\times10^{-5}$ ; $2{,}3\times10^{-10}$ : la fonction renvoie $\res{5}$. Ainsi $1-u_5<10^{-9}$ : $u_5$ est une
    valeur approchée de la limite $1$ à $10^{-9}$ près. Comme $v_{n+1}=v_n^{2}$, le nombre de décimales exactes
    \textbf{double} environ à chaque étape : on parle de convergence quadratique.
\end{questions}
\end{corrige}

%% ---------------------------------------------------------------------
\begin{corrige}{14}{Les abonnés d'une plateforme}
\partieExo{A}{Modélisation}
\begin{questions}
  \item $u_1=0{,}85\times50+12=\res{54{,}5}$\quad et\quad $u_2=0{,}85\times54{,}5+12=\res{58{,}325}$.
    Au 1\textsuperscript{er} janvier 2028, la plateforme compterait environ $58\,325$ abonnés.
  \item Chaque année, 15\,\% des abonnés partent : il en reste 85\,\%, soit $0{,}85\,u_n$ milliers ;
    $12$ milliers de nouveaux abonnés s'y ajoutent. D'où $u_{n+1}=0{,}85\,u_n+12$.
  \item $v_{n+1}=u_{n+1}-80=0{,}85\,u_n-68=0{,}85\,(u_n-80)=0{,}85\,v_n$ : $(v_n)$ est
    $\res{\text{géométrique de raison }0{,}85}$, de premier terme $v_0=50-80=-30$.
  \item $v_n=-30\times0{,}85^{n}$, donc $\res{u_n=80-30\times0{,}85^{n}}$.
  \item $u_{n+1}-u_n=30\times0{,}85^{n}\,(1-0{,}85)=4{,}5\times0{,}85^{n}>0$ : $(u_n)$ est strictement croissante.
    Comme $0<0{,}85<1$, $\lim0{,}85^{n}=0$, donc $\res{\lim u_n=80}$. Le nombre d'abonnés augmente
    et se stabilise vers $80\,000$, sans jamais atteindre ce nombre.
\end{questions}

\noindent
\begin{minipage}[t]{0.52\linewidth}\vspace{0pt}\raggedright
\partieExo{B}{Objectif 75\,000 abonnés}
\begin{questions}[start=6]
  \item On complète avec\\ \lstinline[style=otPython]!while u <= 75:!\\ et \lstinline[style=otPython]!u = 0.85 * u + 12!.
  \item $u_{11}\approx74{,}98\leq75$ et $u_{12}\approx75{,}73>75$ : la fonction renvoie $\res{2038}$.
\end{questions}
\partieExo{C}{Chiffre d'affaires}
\begin{questions}[start=8]
  \item Par somme géométrique :
    \[ S_n=80(n+1)-30\times\frac{1-0{,}85^{n+1}}{1-0{,}85} \]
    $\res{S_n=80(n+1)-200\left(1-0{,}85^{n+1}\right)}$
  \item De 2026 à 2035, $n$ varie de $0$ à $9$ :\\
    $S_9=800-200\left(1-0{,}85^{10}\right)\approx639{,}375$ (milliers).\\
    Chiffre d'affaires : $120\times639\,374{,}9$, soit environ\\
    $\res{76\,725\,000\text{ €}}$.
\end{questions}
\end{minipage}\hfill
\begin{minipage}[t]{0.44\linewidth}\vspace{0pt}
\centering
\begin{tikzpicture}
\begin{axis}[otSchema, width=70mm, height=52mm,
  xmin=0, xmax=21, ymin=40, ymax=86, xtick={0,4,...,20}, ytick={50,60,70,80},
  xlabel={$n$}, xlabel style={at={(1,0)},anchor=west},
  grid=major, grid style={otRule!70,line width=0.3pt}]
  \draw[otNavy,dashed,line width=0.6pt] (axis cs:0,80) -- (axis cs:21,80);
  \draw[otGreen,line width=0.6pt] (axis cs:0,75) -- (axis cs:21,75);
  \addplot[only marks,mark=*,mark size=1.3pt,draw=otNavy,fill=white,samples at={0,...,20}]
    {80-30*0.85^x};
  \addplot[only marks,mark=*,mark size=1.9pt,otGreen] coordinates {(12,75.7327)};
  \node[anchor=north,otEtiquette,font=\scriptsize,text=otGreenDark] at (axis cs:12,74.2) {$u_{12}$};
  \node[anchor=east,font=\scriptsize,text=otGreenDark] at (axis cs:0,75) {$75$};
  \node[anchor=south west,font=\scriptsize,text=otNavy] at (axis cs:14.5,80.4) {$y=80$};
\end{axis}
\end{tikzpicture}\\[0.5mm]
\legende{Nombre d'abonnés (en milliers) : la suite croît vers $80$.}
\end{minipage}
\end{corrige}

%% ---------------------------------------------------------------------
\begin{corrige}{15}{Le problème de Bâle}
\partieExo{A}{Convergence de $(u_n)$}
\begin{questions}
  \item $u_1=\res{1}$,\quad $u_2=1+\frac14=\res{\frac54}$,\quad $u_3=\frac54+\frac19=\res{\frac{49}{36}}$.
  \item $u_{n+1}-u_n=\dfrac{1}{(n+1)^{2}}>0$ : $\res{(u_n)\text{ est croissante}}$.
  \item \begin{questions}
      \item $\dfrac{1}{k}-\dfrac{1}{k+1}=\dfrac{1}{k(k+1)}$, et $k(k+1)\leq(k+1)^{2}$ donne $\dfrac{1}{(k+1)^{2}}\leq\dfrac{1}{k(k+1)}$.
      \item \etape{Initialisation} $u_1=1\leq2-1$.\enspace
        \etape{Hérédité} Si $u_k\leq2-\frac1k$, alors, d'après a),
        $u_{k+1}=u_k+\dfrac{1}{(k+1)^{2}}\leq2-\dfrac1k+\dfrac1k-\dfrac{1}{k+1}=\res{2-\dfrac{1}{k+1}}$.
    \end{questions}
  \item Pour tout $n\geq1$, $u_n\leq2-\frac1n<2$ : la suite $(u_n)$, croissante et majorée par $2$, converge
    vers un réel $\ell$ ; en passant à la limite dans $u_n\leq2$, on obtient $\res{\ell\leq2}$.
\end{questions}

\noindent
\begin{minipage}[t]{0.54\linewidth}\vspace{0pt}\raggedright
\partieExo{B}{Encadrement de $\ell$}
\begin{questions}[start=5]
  \item $w_{n+1}-w_n=\dfrac{1}{(n+1)^{2}}+\dfrac{1}{n+1}-\dfrac{1}{n}$\\[1mm]
    $\phantom{w_{n+1}-w_n}=\dfrac{1}{(n+1)^{2}}-\dfrac{1}{n(n+1)}$\\[1mm]
    $\phantom{w_{n+1}-w_n}=\dfrac{n-(n+1)}{n(n+1)^{2}}=-\dfrac{1}{n(n+1)^{2}}$\\[1mm]
    Ce nombre est négatif : $\res{(w_n)\text{ est décroissante}}$.
  \item $(u_n)$ est croissante et converge vers $\ell$, donc $u_n\leq\ell$. De plus
    $w_n=u_n+\frac1n$ converge aussi vers $\ell$ ; décroissante, elle vérifie $w_n\geq\ell$.
    D'où $\res{u_n\leq\ell\leq w_n}$.
  \item On complète avec \lstinline[style=otPython]!s = s + 1/k**2!. On obtient $u_{1000}\approx1{,}643\,93$ et
    $w_{1000}=u_{1000}+0{,}001$, d'où $\res{1{,}643\,93\leq\ell\leq1{,}644\,94}$.
  \item $\dfrac{\pi^{2}}{6}\approx1{,}644\,934$ appartient bien à cet intervalle.
\end{questions}
\end{minipage}\hfill
\begin{minipage}[t]{0.42\linewidth}\vspace{0pt}
\centering
\begin{tikzpicture}
\begin{axis}[otSchema, width=66mm, height=52mm,
  xmin=0, xmax=13, ymin=0.9, ymax=2.1, xtick={1,4,8,12}, ytick={1,1.5,2},
  yticklabels={$1$,$1{,}5$,$2$},
  xlabel={$n$}, xlabel style={at={(1,0)},anchor=west},
  grid=major, grid style={otRule!70,line width=0.3pt}]
  \draw[otMuted,dashed,line width=0.6pt] (axis cs:0,1.6449) -- (axis cs:13,1.6449);
  \addplot[only marks,mark=*,mark size=1.4pt,otNavy] coordinates {
    (1,1) (2,1.25) (3,1.3611) (4,1.4236) (5,1.4636) (6,1.4914) (7,1.5118)
    (8,1.5274) (9,1.5398) (10,1.5498) (11,1.5580) (12,1.5650)};
  \addplot[only marks,mark=*,mark size=1.4pt,otGreen] coordinates {
    (1,2) (2,1.75) (3,1.6944) (4,1.6736) (5,1.6636) (6,1.6581) (7,1.6547)
    (8,1.6524) (9,1.6509) (10,1.6498) (11,1.6489) (12,1.6483)};
  \node[anchor=east,font=\scriptsize,text=otMuted] at (axis cs:0,1.69) {$\ell$};
  \node[anchor=west,font=\scriptsize,text=otGreenDark] at (axis cs:2.2,1.83) {$w_n$};
  \node[anchor=west,font=\scriptsize,text=otNavy] at (axis cs:2.2,1.17) {$u_n$};
\end{axis}
\end{tikzpicture}\\[0.5mm]
\legende{$(u_n)$ croît et $(w_n)$ décroît : $\ell$ est pris en étau.}
\end{minipage}
\end{corrige}
